\capitulo{4}{Técnicas y herramientas}

El presente capítulo tiene como objetivo describir las técnicas metodológicas y las herramientas de desarrollo que se han utilizado para llevar a cabo el proyecto. Se detallarán los aspectos más destacados de cada elección, justificando el porqué de su utilización frente a posibles alternativas, y cómo estas herramientas han contribuido al logro de los objetivos planteados. Para mejorar la comprensión, las herramientas se han agrupado en subsecciones según su área de aplicación.

\section{Metodología de desarrollo: Agile con SCRUM}
El desarrollo del prototipo se abordó mediante el marco de trabajo \textbf{SCRUM}, estructurando el trabajo en \textit{sprints} de corta duración. Esta metodología ágil fue fundamental para la adaptación continua, permitiendo una redefinición estratégica del alcance (pivote a un MVP híbrido) tras identificar la complejidad del sistema de layout visual. La gestión del proyecto se realizó con \textbf{Zube} \cite{ZubePlatform} para la planificación de tareas y el seguimiento del progreso en colaboración con el tutor, quien adoptó los roles de \textit{Product Owner} y \textit{Scrum Master}.



La adopción de SCRUM no solo ha facilitado la gestión del proyecto, sino que ha sido el pilar que ha permitido una toma de decisiones ágil ante desafíos técnicos, reorientando el esfuerzo hacia la creación de un MVP robusto y funcional.

\section{Herramientas core para el desarrollo de software}

\subsection{Motor de videojuegos: Unity 6 y Lenguaje C\#}
El prototipo se desarrolló sobre \textbf{Unity 6} con \textbf{C\#}, aprovechando su potente motor de renderizado 3D, su sistema de UI unificado (\texttt{Canvas}) y su capacidad de exportación multiplataforma, con un enfoque en dispositivos Android. El IDE principal fue \textbf{Visual Studio}.

\subsection{Modelado 3D y creación de modelos de robot}
El modelo del robot se creó con \textbf{Blender}, tomando como referencia la estética de kits como LEGO WeDo 2.0. Para la programación, se adoptó un \textbf{enfoque híbrido}: la estructura del programa se basa en una adaptación de la librería \textbf{UBlockly}, mientras que los comandos de acción se implementaron a través de un \textbf{lenguaje de dominio específico (DSL) textual}, que incluye un parser (\texttt{CommandParser}) y un intérprete propio.

\subsection{ Programación visual y textual: UBlockly y DSL propio}

\begin{itemize}
	\item \textbf{Base lógica para programación por bloques:} UBlockly proporcionó una estructura de referencia para la representación lógica de bloques (\texttt{Block}) y la gestión de sus conexiones (\texttt{Connection}). Esta base fue utilizada para implementar los bloques estructurales del sistema, como el bloque de evento inicial y los bucles, acelerando el desarrollo de la parte visual.
	\item \textbf{Desarrollo de un lenguaje de dominio específico (DSL) \cite{lenguaje_DSL}: Para los comandos de acción del robot (movimiento y giro), se optó por diseñar e implementar un DSL textual. Esta decisión permitió un control más directo sobre la semántica de las acciones del robot y simplificó la ejecución, eliminando la necesidad de un motor de layout visual complejo para estos comandos. La implementación incluyó:
	\begin{itemize}
		\item Un Parser (\texttt{CommandParser}) que analiza el código escrito por el usuario.
		\item Un Intérprete (\texttt{RobotController}) que ejecuta la secuencia de instrucciones resultante.
	\end{itemize}
	\item \textbf{Implementación híbrida:} La combinación de bloques visuales para la estructura y código textual para las acciones conforma el núcleo híbrido del prototipo. Esta aproximación no solo es una solución pragmática a los desafíos técnicos, sino que también ofrece un modelo pedagógico de transición, donde los usuarios aprenden la estructura del programa visualmente y practican la codificación de comandos de forma escrita.
\end{itemize}

\subsection{ Sistema de gestión de desafíos y guardado de código}
 La persistencia de datos en el prototipo se aborda de dos maneras complementarias:
\begin{itemize}
	\item \textbf{Gestión de desafíos con Scriptable Objects:} La definición de los desafíos (título, descripción, etc.) se gestiona mediante \textbf{Scriptable Objects} (ChallengeData). Este enfoque de Unity permite desacoplar los datos del contenido de la lógica del juego, facilitando la creación y modificación de nuevos desafíos sin necesidad de cambiar el código. Un gestor (ChallengeLoader) se encarga de la persistencia de la selección del usuario entre escenas.
	\item \textbf{Persistencia del programa de usuario (No implementada en el MVP):} Aunque la arquitectura de UBlockly analizada incluía un sistema de serialización a XML, en la versión actual del prototipo no se ha implementado una funcionalidad de guardado/carga del programa creado por el usuario. La lógica de los desafíos está diseñada para ser resuelta en una sola sesión, enfocándose en la resolución de problemas y la ejecución inmediata del código. La implementación de la persistencia se considera una línea de trabajo futura.
\end{itemize}

\subsection{Herramientas adicionales}
La persistencia de los desafíos se gestionó mediante \textbf{Scriptable Objects}. El control de versiones se realizó con \textbf{Git} y \textbf{GitHub}. Se utilizó \textbf{GitHub Copilot} como herramienta de asistencia a la codificación, y la memoria fue redactada en \textbf{LaTeX} con \textbf{TeXstudio}.



